\chapter{Three-Manifolds of Positive Ricci Curvature}

\section*{1. The evolution of curvature under the Ricci flow}

\begin{remark}[Evolution formulas are dimension-independent]
\label{ch06-rem-evolution-formulas-dimension-independent}
\source{chow-knopf-ricci-flow:page-0187}
\dcref{Remark 6.4}
\group{evolution-of-curvature-under-the-ricci-flow}

Unless explicitly stated otherwise, the curvature evolution formulas in this section hold in every dimension \(n \ge 2\).
\end{remark}

\begin{lemma}[General metric variation formulas]
\label{ch06-lem-general-metric-variation}
\source{chow-knopf-ricci-flow:page-0187, chow-knopf-ricci-flow:page-0188}
\dcref{Lemma 6.5}
\group{evolution-of-curvature-under-the-ricci-flow}

If \(g(t)\) is a smooth one-parameter family of metrics on \(M^n\) satisfying
\[
\frac{\partial}{\partial t}g=h,
\]
then the Levi-Civita connection, the \((3,1)\)-Riemann tensor, the Ricci tensor, the scalar curvature, and the volume form evolve by the standard first-variation formulas stated in the source.
\end{lemma}

\begin{proof}
\source{chow-knopf-ricci-flow:page-0187, chow-knopf-ricci-flow:page-0188}
\sketch These are the classical first-variation identities for the Levi-Civita connection, curvature, scalar curvature, and volume form. They follow by differentiating the local coordinate expressions for the connection and curvatures, using the metric-compatibility of \(\nabla\), and tracing the resulting formulas.
\end{proof}

\begin{corollary}[Curvature evolution under Ricci flow]
\label{ch06-cor-curvature-evolution-under-ricci-flow}
\source{chow-knopf-ricci-flow:page-0188}
\dcref{Corollary 6.6}
\group{evolution-of-curvature-under-the-ricci-flow}

If \(g(t)\) is a solution of the Ricci flow, then the Levi-Civita connection, \((3,1)\)-Riemann tensor, Ricci tensor, scalar curvature, and volume form satisfy the evolution equations obtained from Lemma 6.5 after substituting \(h=-2\Rc\).
\end{corollary}

\begin{lemma}[Scalar curvature under Ricci flow]
\label{ch06-lem-scalar-curvature-under-ricci-flow}
\source{chow-knopf-ricci-flow:page-0189}
\dcref{Lemma 6.7}
\group{evolution-of-curvature-under-the-ricci-flow}

The scalar curvature of a solution to the Ricci flow evolves by
\[
\frac{\partial}{\partial t}R=\Delta R+2|\Rc|^2.
\]
\end{lemma}

\begin{lemma}[Preservation of scalar lower bounds]
\label{ch06-lem-preservation-scalar-lower-bounds}
\source{chow-knopf-ricci-flow:page-0189}
\dcref{Lemma 6.8}
\group{evolution-of-curvature-under-the-ricci-flow}

If the scalar curvature of the initial metric \(g_0\) is bounded below by \(\rho\), then the scalar curvature remains bounded below by \(\rho\) for as long as the Ricci flow exists.
\end{lemma}

\begin{lemma}[Ricci evolution in general dimension]
\label{ch06-lem-ricci-evolution-general-dimension}
\source{chow-knopf-ricci-flow:page-0189, chow-knopf-ricci-flow:page-0190}
\dcref{Lemma 6.9}
\group{evolution-of-curvature-under-the-ricci-flow}

Under the Ricci flow, the Ricci tensor evolves by the Lichnerowicz heat equation
\[
\frac{\partial}{\partial t}R_{jk}=\Delta_L R_{jk}=\Delta R_{jk}+2g^{pq}g^{rs}R_{pjkr}R_{qs}-2g^{pq}R_{jp}R_{qk}.
\]
\end{lemma}

\begin{lemma}[Ricci evolution in dimension three]
\label{ch06-lem-ricci-evolution-dimension-three}
\source{chow-knopf-ricci-flow:page-0190}
\dcref{Lemma 6.10}
\group{evolution-of-curvature-under-the-ricci-flow}

In dimension \(n=3\), the Ricci tensor of a solution to the Ricci flow evolves by
\[
\frac{\partial}{\partial t}R_{jk}=\Delta R_{jk}+3RR_{jk}-6g^{pq}R_{jp}R_{qk}+\bigl(2|\Rc|^2-R^2\bigr)g_{jk}.
\]
\end{lemma}

\begin{corollary}[Preservation of Ricci sign in dimension three]
\label{ch06-cor-preservation-ricci-sign-three}
\source{chow-knopf-ricci-flow:page-0190}
\dcref{Corollary 6.11}
\group{evolution-of-curvature-under-the-ricci-flow}

If a three-dimensional solution of the Ricci flow starts with positive, respectively nonnegative, Ricci curvature, then that sign condition is preserved for as long as the solution exists.
\end{corollary}

\begin{remark}[Weyl tensor obstruction]
\label{ch06-rem-weyl-obstruction}
\source{chow-knopf-ricci-flow:page-0190}
\dcref{Remark 6.12}
\group{evolution-of-curvature-under-the-ricci-flow}

In dimensions \(n\ge 3\), the Weyl tensor is the only obstruction to preserving nonnegative Ricci curvature through the evolution equation for \(\Rc\).
\end{remark}

\begin{lemma}[Riemann curvature evolution]
\label{ch06-lem-riemann-curvature-evolution}
\source{chow-knopf-ricci-flow:page-0190, chow-knopf-ricci-flow:page-0191}
\dcref{Lemma 6.13}
\group{evolution-of-curvature-under-the-ricci-flow}

Under the Ricci flow, the \((3,1)\)-Riemann curvature tensor evolves by the stated reaction-diffusion equation with quadratic terms in the curvature tensor and Ricci tensor.
\end{lemma}

\begin{corollary}[Riemann curvature evolution in \((4,0)\)-form]
\label{ch06-cor-riemann-curvature-evolution-four-zero}
\source{chow-knopf-ricci-flow:page-0192}
\dcref{Corollary 6.14}
\group{evolution-of-curvature-under-the-ricci-flow}

Under the Ricci flow, the \((4,0)\)-Riemann curvature tensor satisfies the reaction-diffusion equation obtained by lowering the upper index in Lemma 6.13.
\end{corollary}

\begin{lemma}[Quadratic tensor \(B\)]
\label{ch06-lem-quadratic-tensor-B}
\source{chow-knopf-ricci-flow:page-0192}
\dcref{Equation 6.15}
\group{evolution-of-curvature-under-the-ricci-flow}

The tensor
\[
B_{ijkl} \coloneqq -g^{pr}g^{qs} R_{ipjq}R_{kr\ell s}
\]
is quadratic in the curvature tensor and satisfies the algebraic symmetries \(B_{ijkl}=B_{jilk}=B_{klij}\).
\end{lemma}

\begin{lemma}[Alternative Riemann evolution formula]
\label{ch06-lem-alternative-riemann-evolution}
\source{chow-knopf-ricci-flow:page-0192}
\dcref{Lemma 6.15}
\group{evolution-of-curvature-under-the-ricci-flow}

Under the Ricci flow, the \((4,0)\)-Riemann tensor evolves by the \(B\)-tensor formulation given in the source.
\end{lemma}

\section*{9. Normalized Ricci flow}

\begin{lemma}[Scaling relations]
\label{ch06-lem-scaling-relations}
\source{chow-knopf-ricci-flow:page-0226}
\dcref{Lemma 6.57}
\group{normalized-ricci-flow}

If \(\bar g=\psi g\) for a positive scalar factor \(\psi\), then the Levi-Civita connection and \((3,1)\)-Riemann tensor are unchanged, while the \((4,0)\)-Riemann tensor scales by \(\psi\), the Ricci tensor is unchanged, the scalar curvature scales by \(\psi^{-1}\), and the volume form scales by \(\psi^{n/2}\).
\end{lemma}

\section*{9.2. Long-time existence of the normalized flow}

\begin{theorem}[Long-time existence of the normalized flow]
\label{ch06-thm-long-time-existence-normalized-flow}
\source{chow-knopf-ricci-flow:page-0227, chow-knopf-ricci-flow:page-0228}
\dcref{Theorem 6.58}
\group{normalized-ricci-flow}

If \((\M^3,g(t))\) is the unnormalized Ricci flow starting from a compact three-manifold of positive Ricci curvature, then the corresponding normalized solution exists for all positive time.
\end{theorem}

\begin{lemma}[Blow-up of the maximal scalar curvature integral]
\label{ch06-lem-blowup-max-scalar-integral}
\source{chow-knopf-ricci-flow:page-0227}
\dcref{Lemma 6.59}
\group{normalized-ricci-flow}

For a three-dimensional Ricci flow with initially positive Ricci curvature and maximal existence time \(T<\infty\), one has
\[
\int_0^T R_{\max}(t)\,dt=\infty.
\]
\end{lemma}

\begin{lemma}[Uniform upper bound for the normalized maximum scalar curvature integral]
\label{ch06-lem-uniform-upper-bound-normalized-max-scalar}
\source{chow-knopf-ricci-flow:page-0229}
\dcref{Lemma 6.60}
\group{normalized-ricci-flow}

Along the normalized flow, the time integral of the maximal scalar curvature is bounded above by a constant multiple of the normalized time.
\end{lemma}

\begin{corollary}[Asymptotic Einstein behavior]
\label{ch06-cor-asymptotic-einstein}
\source{chow-knopf-ricci-flow:page-0229}
\dcref{Corollary 6.61}
\group{normalized-ricci-flow}

The normalized flow on a compact three-manifold with initially positive Ricci curvature exists for all positive time and asymptotically approaches an Einstein metric in the sense that the trace-free Ricci tensor becomes uniformly negligible compared to the scalar curvature.
\end{corollary}

\section*{9.3. Evolution of curvature under the normalized flow}

\begin{lemma}[Pure scaling evolution]
\label{ch06-lem-pure-scaling-evolution}
\source{chow-knopf-ricci-flow:page-0230}
\dcref{Lemma 6.62}
\group{normalized-ricci-flow}

If \(\partial_t g=\varphi g\) with \(\varphi\) depending only on time, then the connection and \((3,1)\)-curvature are time-independent, the \((4,0)\)-curvature scales by \(\varphi\), the Ricci tensor is time-independent, and the scalar curvature evolves by \(\partial_t R=-\varphi R\).
\end{lemma}

\begin{remark}[Scaling interpretation]
\label{ch06-rem-scaling-interpretation}
\source{chow-knopf-ricci-flow:page-0230}
\dcref{Remark 6.63}
\group{normalized-ricci-flow}

The formulas of Lemma 6.62 may be deduced from the scaling behavior of the geometric quantities under a time-dependent homothety.
\end{remark}

\begin{corollary}[Curvature evolution under the normalized flow]
\label{ch06-cor-curvature-evolution-normalized-flow}
\source{chow-knopf-ricci-flow:page-0230, chow-knopf-ricci-flow:page-0231}
\dcref{Corollary 6.64}
\group{normalized-ricci-flow}

For a solution of the normalized Ricci flow, the Levi-Civita connection, Riemann curvature, Ricci tensor, and scalar curvature satisfy the normalized evolution equations obtained by combining the unnormalized formulas with the scaling correction \(r(t)\).
\end{corollary}

\section*{10. Exponential convergence}

\begin{lemma}[Pinching of Ricci eigenvalues]
\label{ch06-lem-pinching-ricci-eigenvalues}
\source{chow-knopf-ricci-flow:page-0232}
\dcref{Lemma 6.65}
\group{exponential-convergence}

For a normalized Ricci flow solution on a closed three-manifold of initially positive Ricci curvature, there exists a constant \(B\) such that the largest Ricci eigenvalue \(\lambda\) satisfies
\[
\lambda\le B(\mu+\nu)
\]
for all positive time.
\end{lemma}

\begin{remark}[Scale-invariant pinching form]
\label{ch06-rem-scale-invariant-pinching}
\source{chow-knopf-ricci-flow:page-0232}
\dcref{Remark 6.66}
\group{exponential-convergence}

The estimate of Lemma 6.65 can equivalently be written in the scale-invariant form \(\lambda/(\mu+\nu)\le B\).
\end{remark}

\begin{lemma}[Uniform positive scalar lower bound]
\label{ch06-lem-uniform-positive-scalar-lower-bound}
\source{chow-knopf-ricci-flow:page-0232}
\dcref{Lemma 6.67}
\group{exponential-convergence}

Along the normalized Ricci flow on a closed three-manifold of initially positive Ricci curvature, the scalar curvature is bounded below by a positive constant depending only on the initial metric.
\end{lemma}

\begin{proposition}[Exponential pinching of sectional curvature]
\label{ch06-prop-exponential-pinching-sectional-curvature}
\source{chow-knopf-ricci-flow:page-0232, chow-knopf-ricci-flow:page-0233}
\dcref{Proposition 6.68}
\group{exponential-convergence}

For a normalized Ricci flow solution on a closed three-manifold of initially positive Ricci curvature, there exist constants \(\alpha\in(0,1)\), \(\beta\), and \(C\) depending only on the initial metric such that
\[
\lambda-\nu\le C(\mu+\nu)^{1-\alpha}e^{-\beta t}
\]
for all positive time.
\end{proposition}

\begin{corollary}[Exponential decay of the trace-free Ricci tensor]
\label{ch06-cor-exponential-decay-tracefree-ricci}
\source{chow-knopf-ricci-flow:page-0233}
\dcref{Corollary 6.69}
\group{exponential-convergence}

For a normalized Ricci flow solution on a closed three-manifold of initially positive Ricci curvature, the trace-free Ricci tensor decays exponentially:
\[
\left|\Rc-\frac13 Rg\right| \le Be^{-\beta t}.
\]
\end{corollary}
